\section{Introduction}

An orbit equation records where periodic trajectories live, but it need not
record how their linearized dynamics differ.  This loss is acute when two
periodic points have the same value of a generating action: the action image
identifies the two points, while their return matrices may remain different.
The present paper shows that this apparent information loss produces a
canonical arithmetic decoration in a concrete area-preserving H\'enon
family.

We study
\begin{equation}
\HA(q,p)=(1-Aq^2-p,q),                                      \label{eq:map}
\end{equation}
in the area-preserving quadratic tradition initiated by H\'enon
\citep{henon1969},
with a diagonal reversor-line initial state.  Its period-five critical
points are encoded by a cyclic action \(\PhiA\).  Eliminating the orbit
coordinate gives a plane curve \(\WA(A,c)=0\), where \(c\) is the action
value.  The normalization of this curve is not new: an exact linear
subresultant recovers the familiar six-sheet period-five marker cover.  The
new information appears where the action map fails to separate two points.

The action discriminant contains a degree-nine factor \(\PN(A)^2\) that is
coprime to the ramification factors of the normalization.  A squared
discriminant factor alone does not determine the singularity.  We therefore
work over the generic collision field and reconstruct the normalization
fiber, tangent cone, and two normalization slopes.  The fiber is a separable
quadratic and the slopes differ, proving that the action image has an
ordinary node.

The stability data then descend in a form that the bare action coordinate
does not display.  Let
\[
h_i=\det(I-D\HA^5(z_i))
\]
at the two points above the node.  Branch exchange fixes \(h_1h_2\), and a
common normalization of both Hill values changes the product by a square.
The resulting class in \(K_9^\times/K_9^{\times2}\) is nontrivial, as shown
by its exact rational field norm.  Figure~\ref{fig:descent} summarizes the
construction.

\begin{figure}[t]
\centering
\begin{tikzpicture}[
  node distance=17mm and 22mm,
  every node/.style={align=center},
  >={Latex[length=2.2mm]}
]
\node (q1) {\((A,q_1;h_1)\)};
\node (q2) [below=of q1] {\((A,q_2;h_2)\)};
\node (collision) [right=of q1, yshift=-8.5mm] {action node\\\((A,c_0)\)};
\node (norm) [right=of collision] {descended class\\\([h_1h_2]\in K_9^\times/K_9^{\times2}\)};
\draw[->] (q1) -- (collision) node[midway,above] {same action};
\draw[->] (q2) -- (collision) node[midway,below] {same action};
\draw[->] (collision) -- (norm) node[midway,above] {quadratic norm};
\draw[<->,dashed] (q1) -- (q2) node[midway,left] {branch exchange};
\end{tikzpicture}
\caption{Two distinct points of the period-five normalization map to one
ordinary action node.  Their individual Hill values are branch-labelled,
but their product descends and defines a nontrivial square class.}
\label{fig:descent}
\end{figure}

The contributions are precise and falsifiable:

\begin{enumerate}[leftmargin=*,itemsep=2pt]
\item We derive the action curve from the chronological recurrence and
      separate the known normalization from its singular action embedding.
\item We prove that the generic degree-nine collision is a two-branch
      transverse ordinary node formed by distinct exact-period-five points.
\item We prove the action-Hessian/Hill identity, exclude both \(+1\) and
      \(-1\) multipliers, and descend the symmetric Hill square class.
\item We prove that this class is nonsquare by an exact rational norm and
      prove \(\Gal(P_9/\Q)=S_9\) from three modular cycle types.
\end{enumerate}

Each generic ingredient has prior art, and we make that boundary explicit.
The result supplies no all-period clock, trace formula, analytic determinant,
or self-adjoint operator.  Its value is narrower: it is a theorem-level
example in which chronological stability adds a nontrivial arithmetic layer
to a singular H\'enon action map.
